\subsubsection*{From the curve calculation to arbitrary dimension}

We give the higher-dimensional step independently, so that the passage
from point values to a geometric morphism does not rest on a
fundamental-group calculation. Let $k\subset\overline{\mathbb Q}$ be
contained in $F\mathbb Q^{\mathrm{ab}}$ for some number field $F$.
We use the preceding curve calculation in precisely the following
form: after every finite extension $L/k$, a pinned relative geometric
morphism $C_{\mathrm{et}}\to(\mathbb G_{m,L})_{\mathrm{et}}$, with
$C$ a smooth geometrically connected affine curve having an $L$-point,
has the geometric-point values of a unit of $C$. Its compatible Kummer
class is the class of that unit. The conclusion below is conditional
only on that already proved curve calculation and the arithmetic
inputs explicitly recorded here.

It is enough that the curve calculation initially represents its
Kummer class after a further finite extension. Indeed, suppose the
pointwise unit is $f_C\in\mathcal O(C)^*$ and, over a finite
extension $L'/L$, the Kummer class is $[g_0]$. At every algebraic
point the value of $f_C/g_0$ is divisible in the multiplicative group
of a finite extension of $L'$ and hence is a root of unity, by the
height input used in the curve calculation. A nonconstant rational
function on a geometrically integral curve has cofinite image in
$\mathbb G_m(\overline L)$, so this ratio is constant. Evaluation at
an $L$-point shows that its constant value is itself divisible in
$L'^*$; hence its Kummer class is zero. Thus the two classes agree
after $L'$. At each finite Kummer level, their difference on $C$
lies in the image of $H^1(L,\mu_n)$ by the Hochschild--Serre exact
sequence, since it is zero geometrically. Evaluation at the $L$-point
is a retraction onto this image and is zero on the difference.
Consequently the classes already agree over $L$. This proves exactly
the curve input used below, without requiring a chosen function to
descend merely because its Kummer class descends.

\emph{The arithmetic Picard lemma.}
For every finite extension $L/k$ and every abelian variety $A/L$,
\[
                 \bigcap_{n\geq1}nA(L)=0.                 \tag{CH.1}
\]
Indeed, choose a number field $F_0$ over which $A$ is defined and large
enough that $L\subset M=F_0\mathbb Q^{\mathrm{ab}}$. Theorem~3.1 of
Tsujimura (2020) states that the Kummer map on $A(M)$ is injective
\cite{Tsujimura2020}. Its proof explicitly factors this map as
\[
 A(M)\longrightarrow\varprojlim_n A(M)/nA(M)
       \lhook\joinrel\longrightarrow H^1(G_M,T A),
\]
so its kernel is exactly $\bigcap_n nA(M)$. Thus (CH.1) follows by
inclusion. Equivalently, the finite nonempty division fibres of a point
in $\bigcap_n nA(M)$ have a compatible inverse-limit point: use the
factorial levels and the finite intersection property. There is no
distinction here between unrelated divisions and coherent divisions.

If $Y/L$ is smooth, projective, and geometrically integral, put
$P=\operatorname{Pic}(Y)$ and let $P^0$ be the subgroup whose geometric
classes are algebraically trivial. Then
\[
                 \bigcap_{n\geq1}nP=0,
                 \qquad T P=\varprojlim_n P[n]=0.         \tag{CH.2}
\]
Here $P^0$ injects into the group of $L$-points of the abelian variety
$\operatorname{Pic}^0_{Y/L}$, and
$N=P/P^0$ is a subgroup of the finitely generated geometric
Neron--Severi group; for the latter finiteness statement see
\cite[Theorem~7.3 and Lemma~7.2]{FA_ChengEtAl2021}.
These assertions do not require the Picard scheme
to represent every rational Picard class: the possible Brauer
obstruction only changes the image of $P^0$ inside
$\operatorname{Pic}^0_{Y/L}(L)$. Injectivity of the passage from actual
line bundles follows from Hilbert~90.

For completeness, let $x\in\bigcap_n nP$. Its image in $N$ vanishes,
since $N$ is finitely generated. Let $e$ annihilate the torsion subgroup
of $N$. For every $m\geq1$, choose $y\in P$ with $em y=x$.
Then the image of $y$ in $N$ is torsion, so $z=ey$ lies in $P^0$ and
$mz=x$. Thus $x$ is infinitely divisible in
$\operatorname{Pic}^0_{Y/L}(L)$, and (CH.1) gives $x=0$. Finally, each
coordinate of a compatible torsion system in $T P$ is divisible by
every positive integer, so the first equality in (CH.2) implies the
second. This proves the lemma, including the torsion in $N$.

\emph{Curves with prescribed points and transverse directions.}
We will use the following elementary consequence of Bertini. In a
smooth geometrically integral projective variety over an infinite
characteristic-zero field, a finite set of rational points lies on a
smooth geometrically integral complete-intersection curve. At any
specified point one may require its tangent line to be transverse to
a given smooth divisor. The construction is over the ground field;
one first makes the chosen points rational by a finite extension when
necessary. To check the imposed conditions, take sufficiently high
powers of an ample bundle. Serre vanishing makes evaluation on the
finite first-order neighbourhoods surjective. The linear systems
vanishing at the prescribed points embed the complement of those
points, and a general member is smooth there and geometrically
irreducible. Smoothness and the required independent differentials at
the finitely many prescribed points are nonempty open conditions on
the first jets. Intersect these opens and repeat until dimension one.
This is the usual prescribed-point version of Bertini; the needed
linear-system and irreducibility statements are Stacks, Lemmas
33.47.1, 33.47.3, and 37.32.3 \cite{FA_StacksBertini}.

\emph{Integral boundary residues and the global unit.}
Let $X/k$ be a connected regular affine scheme of finite type with
$x_0\in X(k)$, and let
\[
       \phi:X_{\mathrm{et}}\longrightarrow
                        (\mathbb G_{m,k})_{\mathrm{et}}
\]
be pinned over $B_k$. Write $\phi(a)$ for its value at a geometric
point. Over a finite extension $L/k$, a point $a\in X(L)$ has
$\phi(a)\in L^*$: the relative morphism of closed-point Galois topoi
forces the residue field of the image to embed in $L$. Base change of
$\phi$ to $L$ is obtained by slicing by the finite etale constant-field
object. These observations justify every evaluation used below.

The scheme $X$ is smooth and geometrically integral. In fact, in
characteristic zero regularity implies smoothness, and the algebraic
closure of $k$ in $k(X)$ is a finite field extension contained in all
normal local rings of $X$; evaluation at $x_0$ embeds it in $k$.
Choose a smooth projective compactification $Y$ of $X$ whose boundary
is a divisor with normal crossings, by characteristic-zero embedded
resolution. Concretely, take a projective closure, resolve it while
preserving the smooth open $X$, and principalize its boundary ideal;
the blowups are projective and supported outside $X$. The strong
resolution and monomialization theorem in
\cite[``The result'' and ``Idea of proof'', pp.~1--3]{FA_EncinasHauser2002}
supplies these operations. Write the boundary prime components over
$k$ as $D_1,\ldots,D_s$.
Only smooth points of individual boundary components will be used.

Pulling back the Kummer torsors of the coordinate $t$ gives
\[
       h=(h_n)_n\in\varprojlim_n H^1_{\mathrm{et}}(X,\mu_n).
                                                               \tag{CH.3}
\]
Let $K=k(Y)$. The map
$H^1(X,\mu_n)\to K^*/K^{*n}$ is injective: a generically trivial
finite etale torsor on a normal integral scheme has a section, by
normalization. For every prime divisor $D_i$, the valuation of the
generic Kummer representative of $h_n$ defines a residue
$r_{i,n}\in\mathbb Z/n$. Compatibility gives
$r_i\in\widehat{\mathbb Z}$.

We claim $r_i\in\mathbb Z$. After a finite extension $L/k$, choose
an $L$-rational point $q$ on a geometric component above $D_i$, away
from the other boundary components, and a smooth complete-intersection
curve $\overline C\subset Y_L$ passing through $q$ transversely to
that component and through $x_0$. Its open part
$C=\overline C\cap X_L$ is an affine smooth geometrically integral
curve with an $L$-point. The curve calculation makes $h|_C$ the Kummer
class of an actual unit; its residue at $q$ is therefore an integer.
That residue is exactly $r_i$.

One can check this equality without an inverse-limit purity argument.
At $q$, the regular local ring $R$ is factorial and the only boundary
has one prime equation $z$. Hence
$\operatorname{Pic}(R[1/z])=0$ and
$R[1/z]^*=R^*z^{\mathbb Z}$. For each $n$, the restricted Kummer
class has a representative $u z^j$, with $u\in R^*$, and residue
$j\bmod n$. On the transverse curve $z$ is a local parameter and
$u$ remains a unit, giving the same residue for every $n$. The base
change from $D_i$ to a geometric component has multiplicity one in
characteristic zero. This proves the claim.

Set $D=\sum_i r_iD_i$, now an integral Cartier divisor on $Y$.
Its class lies in $\bigcap_n n\operatorname{Pic}(Y)$. Indeed, choose
a rational function $a_n\in K^*$ representing the generic class of
$h_n$. Its valuations at divisors inside $X$ are multiples of $n$, by
the Kummer description over a discrete valuation ring; at $D_i$ its
valuation is $r_i\bmod n$. Consequently
\[
                    \operatorname{div}(a_n)=D+nE_n
                                                               \tag{CH.4}
\]
for an integral divisor $E_n$ on the regular scheme $Y$.
Thus $[D]=-n[E_n]$. By (CH.2), $[D]=0$, and there is a rational
function $g\in k(Y)^*$ with $\operatorname{div}(g)=D$. In
particular, $g|_X$ is an actual regular unit.

Subtracting $[g]$ from $h$ leaves a compatible class unramified on
$Y$. Explicitly, the divisor of $a_n/g$ is divisible by $n$, so its
associated line bundle with an $n$th-power trivialization defines a
$\mu_n$-torsor on $Y$ extending $h_n-[g]$. Its uniqueness follows
from injectivity into $K^*/K^{*n}$, so these extensions are compatible
as $n$ varies. The Kummer exact sequences on $Y$ and (CH.2) give
\[
  \varprojlim_n H^1(Y,\mu_n)=\widehat{k^*}.
                                                               \tag{CH.5}
\]
To see that no derived-limit term is being suppressed, the maps
$k^*/k^{*mn}\to k^*/k^{*n}$ are surjective; alternatively, since the
compatible Picard-torsion class is zero, each class is represented
uniquely in $k^*/k^{*n}$, and these representatives are automatically
compatible. Here $\Gamma(Y,\mathcal O_Y)^*=k^*$ by geometric
integrality and properness.

Evaluation at $x_0$ identifies the constant class in (CH.5) with the
actual element $c=\phi(x_0)/g(x_0)\in k^*$. Therefore
\[
        f=cg|_X\in\Gamma(X,\mathcal O_X)^*,\qquad
        h=[f],\qquad f(x_0)=\phi(x_0).                 \tag{CH.6}
\]

\emph{Pointwise equality in every dimension.}
Let $a\in X(\overline k)$. Choose a finite extension $L/k$ making
$a$ rational and a smooth geometrically integral curve
$\overline C\subset Y_L$ through both $x_0$ and $a$; put
$C=\overline C\cap X_L$. In dimension one take $C=X_L$ itself.
The curve theorem gives a unit $f_C$ whose values equal those of
$\phi|_C$. Since both $f_C$ and $f|_C$ have Kummer class $h|_C$,
the unit $u=f_C/(f|_C)$ has zero Kummer class for every $n$. The
injection
\[
  \Gamma(C,\mathcal O_C)^*/\Gamma(C,\mathcal O_C)^{*n}
       \lhook\joinrel\longrightarrow H^1(C,\mu_n)
\]
shows that $u$ has an $n$th root among regular units for every $n$.
Its divisor on $\overline C$ is therefore divisible by every positive
integer and is zero. Thus $u\in L^*$ is constant. Evaluation at
$x_0$ gives $u=1$, and consequently
\[
                         \phi(a)=f(a)
                  \quad\hbox{for all }a\in X(\overline k).
                                                               \tag{CH.7}
\]
In dimension zero $X=\operatorname{Spec}k$ and the assertion is
immediate. Notice that this final argument does not infer equality
from a finite collection of Kummer quotients: it uses the integral
curve realization and equality in every Kummer degree. It also does
not require the residual roots of unity to form a chosen compatible
system.

\emph{Recovery of the entire geometric morphism.}
The assertion just proved holds for every connected regular affine
finite-type source with a rational point, and after every finite
extension of the ground field. It is exactly condition (4.18.5) of
Carlson--Haine--Wolf, Theorem~4.18 \cite{CHW2024}. Their equivalence
with (4.18.1) therefore gives the full pinned reconstruction statement
for this field class, with absolute weak normalization as in their
definition. This application addresses all etale objects, not just
the point map of one morphism.

Here is the categorical part of that implication. Two copies of
$\mathbb G_m$ cover $\mathbb A^1$, so the pointwise units just
constructed give pointwise algebraic maps to $\mathbb A^1$. Affine
coordinates and equality on geometric points of a reduced source
then reconstruct the addition and multiplication relations and hence
a map to every affine target. Apply this construction also to each
sliced pinned morphism over a quasicompact separated etale target
object $V$. Its inverse image is again represented by a
quasicompact separated etale source scheme, by CHW Corollary~A.9;
one works on its affine components and, when needed, after a finite
extension supplying a rational point. Uniqueness on geometric points
glues these maps and makes them compatible with all morphisms of
etale objects. Proposition~4.14 then identifies the original inverse
image functor with the functor of the reconstructed scheme map.
Finally, the $h$-descent part of Theorem~4.18 handles arbitrary
finite-type sources. Thus no openness assertion for an unconstructed
birational Galois homomorphism, and no inference from finite set
coefficients alone, is needed.
